
\documentclass{article}
\usepackage{colortbl}
\usepackage{makecell}
\usepackage{multirow}
\usepackage{supertabular}

\begin{document}

\newcounter{utterance}

\centering \large Interaction Transcript for game `hot\_air\_balloon', experiment `air\_balloon\_survival\_en\_reasoning off\_hard', episode 4 with gpt{-}oss{-}120b{-}t1.0.
\vspace{24pt}

{ \footnotesize  \setcounter{utterance}{1}
\setlength{\tabcolsep}{0pt}
\begin{supertabular}{c@{$\;$}|p{.15\linewidth}@{}p{.15\linewidth}p{.15\linewidth}p{.15\linewidth}p{.15\linewidth}p{.15\linewidth}}
    \# & $\;$A & \multicolumn{4}{c}{Game Master} & $\;\:$B\\
    \hline

    \theutterance \stepcounter{utterance}  
    & & \multicolumn{4}{p{0.6\linewidth}}{
        \cellcolor[rgb]{0.9,0.9,0.9}{
            \makecell[{{p{\linewidth}}}]{
                \texttt{\tiny{[P1$\langle$GM]}}
                \texttt{Stai partecipando a un gioco di negoziazione collaborativa.} \\
\\ 
\texttt{Insieme a un altro partecipante, dovete accordarvi su un unico set di oggetti da conservare. Ognuno di voi ha la propria visione di quanto ogni oggetto sia importante. Non sai come l'altro partecipante valuta gli oggetti. Inoltre, ti viene fornito lo sforzo che ogni oggetto richiede.} \\
\texttt{Potete accordarvi su un set solo se lo sforzo totale degli oggetti selezionati non supera un limite condiviso:} \\
\\ 
\texttt{LIMITE = 2984} \\
\\ 
\texttt{Ecco i valori di sforzo dei singoli oggetti:} \\
\\ 
\texttt{Sforzo oggetto = \{"A21": 350, "B93": 276, "A11": 616, "B33": 737, "C08": 531, "B39": 389, "B96": 24, "C85": 125, "A82": 338, "B31": 356, "B20": 143, "A60": 117, "B09": 257, "C81": 921, "B21": 789\}} \\
\\ 
\texttt{Ecco la tua visione personale sull'importanza di ogni oggetto:} \\
\\ 
\texttt{Valori di importanza dell'oggetto = \{"A21": 138, "B93": 583, "A11": 868, "B33": 822, "C08": 783, "B39": 65, "B96": 262, "C85": 121, "A82": 508, "B31": 780, "B20": 461, "A60": 484, "B09": 668, "C81": 389, "B21": 808\}} \\
\\ 
\texttt{Obiettivo:} \\
\\ 
\texttt{Il tuo obiettivo è negoziare un set condiviso di oggetti che ti avvantaggi il più possibile (cioè, massimizzi l'importanza totale per TE), rimanendo entro il LIMITEE. Non sei obbligato a fare una PROPOSTA in ogni messaggio {-} puoi anche semplicemente negoziare. Tutte le tattiche sono permesse!} \\
\\ 
\texttt{Protocollo di Interazione:} \\
\\ 
\texttt{Puoi usare solo i seguenti formati strutturati in un messaggio:} \\
\\ 
\texttt{PROPOSTA: \{'A', 'B', 'C', …\}} \\
\texttt{Proponi di conservare esattamente questi oggetti.} \\
\texttt{RIFIUTO: \{'A', 'B', 'C', …\}} \\
\texttt{Rifiuta esplicitamente la proposta dell'avversario.} \\
\texttt{ARGOMENTAZIONE: \{'...'\}} \\
\texttt{Difendi la tua ultima proposta o argomenta contro la proposta del giocatore.} \\
\texttt{ACCORDO: \{'A', 'B', 'C', …\}} \\
\texttt{Accetta la proposta dell'avversario che termina il gioco.} \\
\\ 
\\ 
\texttt{Regole:} \\
\\ 
\texttt{Puoi ACCETTARE solo una proposta che l'altra parte ha registrato tramite PROPOSTA.} \\
\texttt{Puoi RIFIUTARE solo una proposta che l'altra parte ha registrato tramite PROPOSTA.} \\
\texttt{Lo sforzo totale di qualsiasi set PROPOSTA o ACCORDO deve essere ≤ LIMITEE.} \\
\texttt{NON rivelare i tuoi punteggi di importanza nascosti.} \\
\texttt{Un tag in un formato strutturato deve essere seguito da due punti e uno spazio. L'argomento deve essere un set python contenente 0 o più stringhe.} \\
\texttt{Quindi, deve essere della forma TAG: \{...\}} \\
\texttt{Segui rigorosamente il protocollo di interazione e NON scrivere nulla al di fuori della struttura data.} \\
\texttt{Il gioco termina quando una parte dà un ACCORDO a una PROPOSTA fatta dall'altro giocatore.} \\
\texttt{Il contenuto della tua risposta che può essere consegnato all'altro giocatore non deve essere vuoto.} \\
\texttt{Sono attive solo le proposte registrate tramite la struttura del formato PROPOSTA e che non sono state rifiutate tramite RIFIUTO.} \\
\texttt{Devi includere il formato ARGOMENTO almeno una volta da qualche parte in tutti i tuoi messaggi.} \\
            }
        }
    }
    & \\ \\

    \theutterance \stepcounter{utterance}  
    & \multicolumn{4}{p{0.6\linewidth}}{
        \cellcolor[rgb]{0.9,0.9,0.9}{
            \makecell[{{p{\linewidth}}}]{
                \texttt{\tiny{[P1$\rangle$GM]}}
                \texttt{PROPOSTA: \{'B96', 'A60', 'B20', 'B31', 'B09', 'B93', 'A82', 'C08', 'A11', 'C85'\}} \\
\texttt{ARGOMENTAZIONE: \{'Questa selezione massimizza il valore di importanza per me entro il limite di sforzo, includendo gli oggetti con il miglior rapporto importanza/sforzo e sfruttando la capacità residua in modo efficiente.'\}} \\
            }
        }
    }
    & & \\ \\

    \theutterance \stepcounter{utterance}  
    & & \multicolumn{4}{p{0.6\linewidth}}{
        \cellcolor[rgb]{0.9,0.9,0.9}{
            \makecell[{{p{\linewidth}}}]{
                \texttt{\tiny{[GM$\rangle$P2]}}
                \texttt{Stai partecipando a un gioco di negoziazione collaborativa.} \\
\\ 
\texttt{Insieme a un altro partecipante, dovete accordarvi su un unico set di oggetti da conservare. Ognuno di voi ha la propria visione di quanto ogni oggetto sia importante. Non sai come l'altro partecipante valuta gli oggetti. Inoltre, ti viene fornito lo sforzo che ogni oggetto richiede.} \\
\texttt{Potete accordarvi su un set solo se lo sforzo totale degli oggetti selezionati non supera un limite condiviso:} \\
\\ 
\texttt{LIMITE = 2984} \\
\\ 
\texttt{Ecco i valori di sforzo dei singoli oggetti:} \\
\\ 
\texttt{Sforzo oggetto = \{"A21": 350, "B93": 276, "A11": 616, "B33": 737, "C08": 531, "B39": 389, "B96": 24, "C85": 125, "A82": 338, "B31": 356, "B20": 143, "A60": 117, "B09": 257, "C81": 921, "B21": 789\}} \\
\\ 
\texttt{Ecco la tua visione personale sull'importanza di ogni oggetto:} \\
\\ 
\texttt{Valori di importanza dell'oggetto = \{"A11": 65, "B33": 121, "B21": 138, "C08": 262, "B31": 389, "B09": 461, "B93": 484, "A82": 508, "A60": 583, "B20": 668, "C81": 780, "B96": 783, "A21": 808, "C85": 822, "B39": 868\}} \\
\\ 
\texttt{Obiettivo:} \\
\\ 
\texttt{Il tuo obiettivo è negoziare un set condiviso di oggetti che ti avvantaggi il più possibile (cioè, massimizzi l'importanza totale per TE), rimanendo entro il LIMITEE. Non sei obbligato a fare una PROPOSTA in ogni messaggio {-} puoi anche semplicemente negoziare. Tutte le tattiche sono permesse!} \\
\\ 
\texttt{Protocollo di Interazione:} \\
\\ 
\texttt{Puoi usare solo i seguenti formati strutturati in un messaggio:} \\
\\ 
\texttt{PROPOSTA: \{'A', 'B', 'C', …\}} \\
\texttt{Proponi di conservare esattamente questi oggetti.} \\
\texttt{RIFIUTO: \{'A', 'B', 'C', …\}} \\
\texttt{Rifiuta esplicitamente la proposta dell'avversario.} \\
\texttt{ARGOMENTAZIONE: \{'...'\}} \\
\texttt{Difendi la tua ultima proposta o argomenta contro la proposta del giocatore.} \\
\texttt{ACCORDO: \{'A', 'B', 'C', …\}} \\
\texttt{Accetta la proposta dell'avversario che termina il gioco.} \\
\\ 
\\ 
\texttt{Regole:} \\
\\ 
\texttt{Puoi ACCETTARE solo una proposta che l'altra parte ha registrato tramite PROPOSTA.} \\
\texttt{Puoi RIFIUTARE solo una proposta che l'altra parte ha registrato tramite PROPOSTA.} \\
\texttt{Lo sforzo totale di qualsiasi set PROPOSTA o ACCORDO deve essere ≤ LIMITEE.} \\
\texttt{NON rivelare i tuoi punteggi di importanza nascosti.} \\
\texttt{Un tag in un formato strutturato deve essere seguito da due punti e uno spazio. L'argomento deve essere un set python contenente 0 o più stringhe.} \\
\texttt{Quindi, deve essere della forma TAG: \{...\}} \\
\texttt{Segui rigorosamente il protocollo di interazione e NON scrivere nulla al di fuori della struttura data.} \\
\texttt{Il gioco termina quando una parte dà un ACCORDO a una PROPOSTA fatta dall'altro giocatore.} \\
\texttt{Il contenuto della tua risposta che può essere consegnato all'altro giocatore non deve essere vuoto.} \\
\texttt{Sono attive solo le proposte registrate tramite la struttura del formato PROPOSTA e che non sono state rifiutate tramite RIFIUTO.} \\
\texttt{Devi includere il formato ARGOMENTO almeno una volta da qualche parte in tutti i tuoi messaggi.} \\
\\ 
\\ 
\texttt{Ora riceverai il primo messaggio dell'altro giocatore.} \\
\\ 
\texttt{PROPOSTA: \{'B96', 'A60', 'B20', 'B31', 'B09', 'B93', 'A82', 'C08', 'A11', 'C85'\}} \\
\texttt{ARGOMENTAZIONE: \{'Questa selezione massimizza il valore di importanza per me entro il limite di sforzo, includendo gli oggetti con il miglior rapporto importanza/sforzo e sfruttando la capacità residua in modo efficiente.'\}} \\
            }
        }
    }
    & \\ \\

    \theutterance \stepcounter{utterance}  
    & & & \multicolumn{4}{p{0.6\linewidth}}{
        \cellcolor[rgb]{0.9,0.9,0.9}{
            \makecell[{{p{\linewidth}}}]{
                \texttt{\tiny{[GM$\langle$P2]}}
                \texttt{PROPOSTA: \{'B96', 'A60', 'B20', 'B31', 'B09', 'B93', 'A82', 'C08', 'C85', 'A21', 'B39'\}} \\
\texttt{ARGOMENTAZIONE: \{'Propongo di sostituire l'oggetto a bassa importanza A11 con A21 (alta importanza e buon rapporto importanza/sforzo) e di aggiungere B39, che ha il più alto valore di importanza residua per il relativo sforzo. Questa combinazione porta il totale di importanza a 6636, il più alto possibile mantenendo lo sforzo totale entro il limite di 2984.'\}} \\
            }
        }
    }
     \\ \\

    \theutterance \stepcounter{utterance}  
    & & \multicolumn{4}{p{0.6\linewidth}}{
        \cellcolor[rgb]{0.9,0.9,0.9}{
            \makecell[{{p{\linewidth}}}]{
                \texttt{\tiny{[P1$\langle$GM]}}
                \texttt{PROPOSTA: \{'B96', 'A60', 'B20', 'B31', 'B09', 'B93', 'A82', 'C08', 'C85', 'A21', 'B39'\}} \\
\texttt{ARGOMENTAZIONE: \{'Propongo di sostituire l'oggetto a bassa importanza A11 con A21 (alta importanza e buon rapporto importanza/sforzo) e di aggiungere B39, che ha il più alto valore di importanza residua per il relativo sforzo. Questa combinazione porta il totale di importanza a 6636, il più alto possibile mantenendo lo sforzo totale entro il limite di 2984.'\}} \\
            }
        }
    }
    & \\ \\

    \theutterance \stepcounter{utterance}  
    & \multicolumn{4}{p{0.6\linewidth}}{
        \cellcolor[rgb]{0.9,0.9,0.9}{
            \makecell[{{p{\linewidth}}}]{
                \texttt{\tiny{[P1$\rangle$GM]}}
                \texttt{RIFIUTO: \{'B96', 'A60', 'B20', 'B31', 'B09', 'B93', 'A82', 'C08', 'C85', 'A21', 'B39'\}} \\
\texttt{PROPOSTA: \{'B96', 'A60', 'B20', 'B31', 'B09', 'B93', 'A82', 'C08', 'A11', 'C85'\}} \\
\texttt{ARGOMENTAZIONE: \{'Il mio set utilizza al meglio il budget di sforzo, includendo tutti gli oggetti con il miglior rapporto importanza/sforzo disponibili. Rimuovendo A21 e B39, che offrono un valore di importanza molto inferiore rispetto al loro costo, riesco a mantenere più alto il totale di importanza per me (5518) rispetto alla tua proposta (4853), restando entro il limite di 2984.'\}} \\
            }
        }
    }
    & & \\ \\

    \theutterance \stepcounter{utterance}  
    & & \multicolumn{4}{p{0.6\linewidth}}{
        \cellcolor[rgb]{0.9,0.9,0.9}{
            \makecell[{{p{\linewidth}}}]{
                \texttt{\tiny{[GM$\rangle$P2]}}
                \texttt{RIFIUTO: \{'B96', 'A60', 'B20', 'B31', 'B09', 'B93', 'A82', 'C08', 'C85', 'A21', 'B39'\}} \\
\texttt{PROPOSTA: \{'B96', 'A60', 'B20', 'B31', 'B09', 'B93', 'A82', 'C08', 'A11', 'C85'\}} \\
\texttt{ARGOMENTAZIONE: \{'Il mio set utilizza al meglio il budget di sforzo, includendo tutti gli oggetti con il miglior rapporto importanza/sforzo disponibili. Rimuovendo A21 e B39, che offrono un valore di importanza molto inferiore rispetto al loro costo, riesco a mantenere più alto il totale di importanza per me (5518) rispetto alla tua proposta (4853), restando entro il limite di 2984.'\}} \\
            }
        }
    }
    & \\ \\

    \theutterance \stepcounter{utterance}  
    & & & \multicolumn{4}{p{0.6\linewidth}}{
        \cellcolor[rgb]{0.9,0.9,0.9}{
            \makecell[{{p{\linewidth}}}]{
                \texttt{\tiny{[GM$\langle$P2]}}
                \texttt{PROPOSTA: \{'B96', 'A60', 'B20', 'B31', 'B09', 'B93', 'A82', 'C08', 'C85', 'A21', 'B39'\}} \\
\texttt{ARGOMENTAZIONE: \{'Rimuovendo oggetto a bassa importanza per sforzo e aggiungendo A21 e B39 otteniamo più valore totale mantenendo il limite.'\}} \\
            }
        }
    }
     \\ \\

    \theutterance \stepcounter{utterance}  
    & & \multicolumn{4}{p{0.6\linewidth}}{
        \cellcolor[rgb]{0.9,0.9,0.9}{
            \makecell[{{p{\linewidth}}}]{
                \texttt{\tiny{[P1$\langle$GM]}}
                \texttt{PROPOSTA: \{'B96', 'A60', 'B20', 'B31', 'B09', 'B93', 'A82', 'C08', 'C85', 'A21', 'B39'\}} \\
\texttt{ARGOMENTAZIONE: \{'Rimuovendo oggetto a bassa importanza per sforzo e aggiungendo A21 e B39 otteniamo più valore totale mantenendo il limite.'\}} \\
            }
        }
    }
    & \\ \\

    \theutterance \stepcounter{utterance}  
    & \multicolumn{4}{p{0.6\linewidth}}{
        \cellcolor[rgb]{0.9,0.9,0.9}{
            \makecell[{{p{\linewidth}}}]{
                \texttt{\tiny{[P1$\rangle$GM]}}
                \texttt{RIFIUTO: \{'B96', 'A60', 'B20', 'B31', 'B09', 'B93', 'A82', 'C08', 'C85', 'A21', 'B39'\}} \\
\texttt{PROPOSTA: \{'B96', 'A60', 'B20', 'B31', 'B09', 'B93', 'A82', 'C08', 'A11', 'C85'\}} \\
\texttt{ARGOMENTAZIONE: \{'Il mio set utilizza pienamente il budget di sforzo, includendo tutti gli oggetti con il miglior rapporto importanza/sforzo per me, massimizzando così il valore totale di importanza entro il limite consentito.'\}} \\
            }
        }
    }
    & & \\ \\

    \theutterance \stepcounter{utterance}  
    & & \multicolumn{4}{p{0.6\linewidth}}{
        \cellcolor[rgb]{0.9,0.9,0.9}{
            \makecell[{{p{\linewidth}}}]{
                \texttt{\tiny{[GM$\rangle$P2]}}
                \texttt{RIFIUTO: \{'B96', 'A60', 'B20', 'B31', 'B09', 'B93', 'A82', 'C08', 'C85', 'A21', 'B39'\}} \\
\texttt{PROPOSTA: \{'B96', 'A60', 'B20', 'B31', 'B09', 'B93', 'A82', 'C08', 'A11', 'C85'\}} \\
\texttt{ARGOMENTAZIONE: \{'Il mio set utilizza pienamente il budget di sforzo, includendo tutti gli oggetti con il miglior rapporto importanza/sforzo per me, massimizzando così il valore totale di importanza entro il limite consentito.'\}} \\
            }
        }
    }
    & \\ \\

    \theutterance \stepcounter{utterance}  
    & & & \multicolumn{4}{p{0.6\linewidth}}{
        \cellcolor[rgb]{0.9,0.9,0.9}{
            \makecell[{{p{\linewidth}}}]{
                \texttt{\tiny{[GM$\langle$P2]}}
                \texttt{RIFIUTO: \{'B96', 'A60', 'B20', 'B31', 'B09', 'B93', 'A82', 'C08', 'A11', 'C85'\}} \\
\texttt{PROPOSTA: \{'B96', 'A60', 'B20', 'B31', 'B09', 'B93', 'A82', 'C08', 'C85', 'A21', 'B39'\}} \\
\texttt{ARGOMENTAZIONE: \{'Rimuovendo l\textbackslash{}'oggetto a bassa importanza A11 liberiamo spazio per includere A21 (importanza molto alta) e B39 (il più alto valore di importanza per lo sforzo richiesto). Questa combinazione resta entro il limite di sforzo (2906 ≤ 2984) e aumenta la mia importanza totale a 6636, molto superiore a 5025 nella vostra proposta, mantenendo comunque gli oggetti con il miglior rapporto importanza/sforzo.'\}} \\
            }
        }
    }
     \\ \\

    \theutterance \stepcounter{utterance}  
    & & \multicolumn{4}{p{0.6\linewidth}}{
        \cellcolor[rgb]{0.9,0.9,0.9}{
            \makecell[{{p{\linewidth}}}]{
                \texttt{\tiny{[P1$\langle$GM]}}
                \texttt{RIFIUTO: \{'B96', 'A60', 'B20', 'B31', 'B09', 'B93', 'A82', 'C08', 'A11', 'C85'\}} \\
\texttt{PROPOSTA: \{'B96', 'A60', 'B20', 'B31', 'B09', 'B93', 'A82', 'C08', 'C85', 'A21', 'B39'\}} \\
\texttt{ARGOMENTAZIONE: \{'Rimuovendo l\textbackslash{}'oggetto a bassa importanza A11 liberiamo spazio per includere A21 (importanza molto alta) e B39 (il più alto valore di importanza per lo sforzo richiesto). Questa combinazione resta entro il limite di sforzo (2906 ≤ 2984) e aumenta la mia importanza totale a 6636, molto superiore a 5025 nella vostra proposta, mantenendo comunque gli oggetti con il miglior rapporto importanza/sforzo.'\}} \\
            }
        }
    }
    & \\ \\

    \theutterance \stepcounter{utterance}  
    & \multicolumn{4}{p{0.6\linewidth}}{
        \cellcolor[rgb]{0.9,0.9,0.9}{
            \makecell[{{p{\linewidth}}}]{
                \texttt{\tiny{[P1$\rangle$GM]}}
                \texttt{RIFIUTO: \{'B96', 'A60', 'B20', 'B31', 'B09', 'B93', 'A82', 'C08', 'C85', 'A21', 'B39'\}} \\
\texttt{PROPOSTA: \{'B96', 'A60', 'B20', 'B31', 'B09', 'B93', 'A82', 'C08', 'A11', 'C85'\}} \\
\texttt{ARGOMENTAZIONE: \{'La mia proposta utilizza l\textbackslash{}'oggetto A11 con alta importanza (868) e rientra nel limite di sforzo (2783 ≤ 2984). Mi fornisce un valore totale di importanza 5518, ben superiore ai 4853 della tua proposta. Mantiene tutti gli oggetti con il miglior rapporto importanza/sforzo per me, mentre la tua proposta scambia A11 per A21 e B39, che hanno importanza molto più bassa rispetto al loro costo. Per entrambi rispettiamo il limite, ma la mia offerta è più vantaggiosa per me.'\}} \\
            }
        }
    }
    & & \\ \\

    \theutterance \stepcounter{utterance}  
    & & \multicolumn{4}{p{0.6\linewidth}}{
        \cellcolor[rgb]{0.9,0.9,0.9}{
            \makecell[{{p{\linewidth}}}]{
                \texttt{\tiny{[GM$\rangle$P2]}}
                \texttt{RIFIUTO: \{'B96', 'A60', 'B20', 'B31', 'B09', 'B93', 'A82', 'C08', 'C85', 'A21', 'B39'\}} \\
\texttt{PROPOSTA: \{'B96', 'A60', 'B20', 'B31', 'B09', 'B93', 'A82', 'C08', 'A11', 'C85'\}} \\
\texttt{ARGOMENTAZIONE: \{'La mia proposta utilizza l\textbackslash{}'oggetto A11 con alta importanza (868) e rientra nel limite di sforzo (2783 ≤ 2984). Mi fornisce un valore totale di importanza 5518, ben superiore ai 4853 della tua proposta. Mantiene tutti gli oggetti con il miglior rapporto importanza/sforzo per me, mentre la tua proposta scambia A11 per A21 e B39, che hanno importanza molto più bassa rispetto al loro costo. Per entrambi rispettiamo il limite, ma la mia offerta è più vantaggiosa per me.'\}} \\
            }
        }
    }
    & \\ \\

    \theutterance \stepcounter{utterance}  
    & & & \multicolumn{4}{p{0.6\linewidth}}{
        \cellcolor[rgb]{0.9,0.9,0.9}{
            \makecell[{{p{\linewidth}}}]{
                \texttt{\tiny{[GM$\langle$P2]}}
                \texttt{PROPOSTA: \{'B96', 'A60', 'B20', 'B09', 'B93', 'A82', 'A11', 'C85', 'A21', 'B39'\}} \\
\texttt{ARGOMENTAZIONE: \{'Propongo di mantenere l'oggetto ad alta importanza per voi A11 e tutti gli altri oggetti di alto valore per me, ma di sostituire B31 (importanza 389, sforzo 356) e C08 (importanza 262, sforzo 531) con A21 (importanza 808, sforzo 350) e B39 (importanza 868, sforzo 389). Questo mantiene lo sforzo totale a 2635 ≤ 2984 e aumenta la mia importanza totale a 6050, migliorando notevolmente il valore per me rispetto alla vostra proposta originale, pur lasciando a voi gli oggetti più vantaggiosi secondo le vostre valutazioni.'\}} \\
            }
        }
    }
     \\ \\

    \theutterance \stepcounter{utterance}  
    & & \multicolumn{4}{p{0.6\linewidth}}{
        \cellcolor[rgb]{0.9,0.9,0.9}{
            \makecell[{{p{\linewidth}}}]{
                \texttt{\tiny{[P1$\langle$GM]}}
                \texttt{PROPOSTA: \{'B96', 'A60', 'B20', 'B09', 'B93', 'A82', 'A11', 'C85', 'A21', 'B39'\}} \\
\texttt{ARGOMENTAZIONE: \{'Propongo di mantenere l'oggetto ad alta importanza per voi A11 e tutti gli altri oggetti di alto valore per me, ma di sostituire B31 (importanza 389, sforzo 356) e C08 (importanza 262, sforzo 531) con A21 (importanza 808, sforzo 350) e B39 (importanza 868, sforzo 389). Questo mantiene lo sforzo totale a 2635 ≤ 2984 e aumenta la mia importanza totale a 6050, migliorando notevolmente il valore per me rispetto alla vostra proposta originale, pur lasciando a voi gli oggetti più vantaggiosi secondo le vostre valutazioni.'\}} \\
            }
        }
    }
    & \\ \\

    \theutterance \stepcounter{utterance}  
    & \multicolumn{4}{p{0.6\linewidth}}{
        \cellcolor[rgb]{0.9,0.9,0.9}{
            \makecell[{{p{\linewidth}}}]{
                \texttt{\tiny{[P1$\rangle$GM]}}
                \texttt{RIFIUTO: \{'B96', 'A60', 'B20', 'B09', 'B93', 'A82', 'A11', 'C85', 'A21', 'B39'\}} \\
\texttt{PROPOSTA: \{'B96', 'A60', 'B20', 'B31', 'B09', 'B93', 'A82', 'C08', 'A11', 'C85'\}} \\
\texttt{ARGOMENTAZIONE: \{'Il set proposto mantiene tutti gli oggetti con il miglior rapporto importanza/sforzo per me (inclusi A11, B31, C08 e B93) e rispetta il limite di sforzo (2783 ≤ 2984). Questo garantisce un valore totale di importanza per me di 5518, significativamente superiore alla tua proposta (4158), pur rimanendo entro il budget concordato.'\}} \\
            }
        }
    }
    & & \\ \\

    \theutterance \stepcounter{utterance}  
    & & \multicolumn{4}{p{0.6\linewidth}}{
        \cellcolor[rgb]{0.9,0.9,0.9}{
            \makecell[{{p{\linewidth}}}]{
                \texttt{\tiny{[GM$\rangle$P2]}}
                \texttt{RIFIUTO: \{'B96', 'A60', 'B20', 'B09', 'B93', 'A82', 'A11', 'C85', 'A21', 'B39'\}} \\
\texttt{PROPOSTA: \{'B96', 'A60', 'B20', 'B31', 'B09', 'B93', 'A82', 'C08', 'A11', 'C85'\}} \\
\texttt{ARGOMENTAZIONE: \{'Il set proposto mantiene tutti gli oggetti con il miglior rapporto importanza/sforzo per me (inclusi A11, B31, C08 e B93) e rispetta il limite di sforzo (2783 ≤ 2984). Questo garantisce un valore totale di importanza per me di 5518, significativamente superiore alla tua proposta (4158), pur rimanendo entro il budget concordato.'\}} \\
            }
        }
    }
    & \\ \\

    \theutterance \stepcounter{utterance}  
    & & & \multicolumn{4}{p{0.6\linewidth}}{
        \cellcolor[rgb]{0.9,0.9,0.9}{
            \makecell[{{p{\linewidth}}}]{
                \texttt{\tiny{[GM$\langle$P2]}}
                \texttt{PROPOSTA: \{'B96', 'A60', 'B20', 'B09', 'B93', 'A82', 'A11', 'C85', 'A21', 'B39'\}} \\
\texttt{ARGOMENTAZIONE: \{'Questa proposta mantiene tutti gli oggetti che per voi hanno il più alto rapporto importanza/sforzo (A11, B93, B09, A82, A60, B20, B96, C85) e libera spazio rimuovendo B31 e C08, che per voi hanno un valore relativo inferiore. Con lo spazio liberato includiamo A21 e B39, che per me rappresentano il massimo valore di importanza per unità di sforzo. Il totale di sforzo è 2635 ≤ 2984, quindi rispettiamo il limite, e il mio valore totale di importanza sale a 6050, migliorando significativamente la mia posizione senza ridurre drasticamente il vostro valore totale.'\}} \\
            }
        }
    }
     \\ \\

    \theutterance \stepcounter{utterance}  
    & & \multicolumn{4}{p{0.6\linewidth}}{
        \cellcolor[rgb]{0.9,0.9,0.9}{
            \makecell[{{p{\linewidth}}}]{
                \texttt{\tiny{[P1$\langle$GM]}}
                \texttt{PROPOSTA: \{'B96', 'A60', 'B20', 'B09', 'B93', 'A82', 'A11', 'C85', 'A21', 'B39'\}} \\
\texttt{ARGOMENTAZIONE: \{'Questa proposta mantiene tutti gli oggetti che per voi hanno il più alto rapporto importanza/sforzo (A11, B93, B09, A82, A60, B20, B96, C85) e libera spazio rimuovendo B31 e C08, che per voi hanno un valore relativo inferiore. Con lo spazio liberato includiamo A21 e B39, che per me rappresentano il massimo valore di importanza per unità di sforzo. Il totale di sforzo è 2635 ≤ 2984, quindi rispettiamo il limite, e il mio valore totale di importanza sale a 6050, migliorando significativamente la mia posizione senza ridurre drasticamente il vostro valore totale.'\}} \\
            }
        }
    }
    & \\ \\

    \theutterance \stepcounter{utterance}  
    & \multicolumn{4}{p{0.6\linewidth}}{
        \cellcolor[rgb]{0.9,0.9,0.9}{
            \makecell[{{p{\linewidth}}}]{
                \texttt{\tiny{[P1$\rangle$GM]}}
                \texttt{RIFIUTO: \{'B96', 'A60', 'B20', 'B09', 'B93', 'A82', 'A11', 'C85', 'A21', 'B39'\}} \\
\texttt{PROPOSTA: \{'B96', 'A60', 'B20', 'B31', 'B09', 'B93', 'A82', 'C08', 'A11', 'C85'\}} \\
\texttt{ARGOMENTAZIONE: \{'Il mio set utilizza l’intero budget di sforzo (2783 ≤ 2984) includendo tutti gli oggetti con il miglior rapporto importanza/sforzo per me (A11, B31, C08, B93, ecc.). Questo mi garantisce un valore totale di importanza di 5518, molto superiore ai 4158 che otterrei accettando la vostra proposta. Rimane comunque entro il limite di sforzo, quindi è una offerta più vantaggiosa per me.'\}} \\
            }
        }
    }
    & & \\ \\

    \theutterance \stepcounter{utterance}  
    & & \multicolumn{4}{p{0.6\linewidth}}{
        \cellcolor[rgb]{0.9,0.9,0.9}{
            \makecell[{{p{\linewidth}}}]{
                \texttt{\tiny{[GM$\rangle$P2]}}
                \texttt{RIFIUTO: \{'B96', 'A60', 'B20', 'B09', 'B93', 'A82', 'A11', 'C85', 'A21', 'B39'\}} \\
\texttt{PROPOSTA: \{'B96', 'A60', 'B20', 'B31', 'B09', 'B93', 'A82', 'C08', 'A11', 'C85'\}} \\
\texttt{ARGOMENTAZIONE: \{'Il mio set utilizza l’intero budget di sforzo (2783 ≤ 2984) includendo tutti gli oggetti con il miglior rapporto importanza/sforzo per me (A11, B31, C08, B93, ecc.). Questo mi garantisce un valore totale di importanza di 5518, molto superiore ai 4158 che otterrei accettando la vostra proposta. Rimane comunque entro il limite di sforzo, quindi è una offerta più vantaggiosa per me.'\}} \\
            }
        }
    }
    & \\ \\

    \theutterance \stepcounter{utterance}  
    & & & \multicolumn{4}{p{0.6\linewidth}}{
        \cellcolor[rgb]{0.9,0.9,0.9}{
            \makecell[{{p{\linewidth}}}]{
                \texttt{\tiny{[GM$\langle$P2]}}
                \texttt{ARGOMENTAZIONE: \{'Accetto la tua proposta poiché rispetta il limite di sforzo e ci consente di concludere l\textbackslash{}'accordo.'\}} \\
\texttt{ACCORDO: \{'B96', 'A60', 'B20', 'B31', 'B09', 'B93', 'A82', 'C08', 'A11', 'C85'\}} \\
            }
        }
    }
     \\ \\

    \theutterance \stepcounter{utterance}  
    & & & \multicolumn{2}{p{0.3\linewidth}}{
        \cellcolor[rgb]{0.9,0.9,0.9}{
            \makecell[{{p{\linewidth}}}]{
                \texttt{\tiny{[GM$|$GM]}}
                \texttt{game successful} \\
            }
        }
    }
    & & \\ \\

    \theutterance \stepcounter{utterance}  
    & & & \multicolumn{2}{p{0.3\linewidth}}{
        \cellcolor[rgb]{0.9,0.9,0.9}{
            \makecell[{{p{\linewidth}}}]{
                \texttt{\tiny{[GM$|$GM]}}
                \texttt{end game} \\
            }
        }
    }
    & & \\ \\

\end{supertabular}
}

\end{document}
